% ID: SEL-D0034-OA-E04
% ESTADO: cotejada. Segunda pasada visual y borrador compilado revisados.
% Fuente: Selectividad/matematicas_aplicadas/titular_septiembre_1_MATEMATICAS_APLICADAS.pdf
% SHA-256: c637bbde5c26c0b35b59202bd832ddc97990d3845365ce483fbf2a49832f73fe
% Página: 1; opción A; ejercicio original 4.

Una cadena de hipermercados decide estudiar la proporción de artículos de un determinado tipo que tienen defectos en su envoltorio. Para ello, selecciona aleatoriamente 2000 artículos de este tipo entre sus hipermercados y encuentra que 19 de ellos tienen defectos en su envoltorio.
\begin{enumerate}[label=\alph*)]
\item \textbf{(1.5 puntos)} Determine un intervalo, al 95\% de confianza, para la proporción real de artículos con este tipo de defecto e interprete el resultado obtenido.
\item \textbf{(1 punto)} \textquestiondown{}Cuántos artículos, como mínimo, deberá seleccionar para que, con un nivel de confianza del 99\%, la proporción muestral difiera de la proporción real a lo sumo en un 1\%?
\end{enumerate}
